% GENERATED FILE. Edit canonical lessons, then run npm run generate:books.
% canonical-source-hash: fnv1a64:af93fec238f158c9
% canonical-lessons: ML-C392-anusocanam, ML-C392-adaranjali, ML-C392-caramavarsikam, ML-C392-semitteri, ML-C392-anusocikkuka

\chapter{Mourning}
\label{ch:ml-mourning}
\hlchaptermodality{392}

\begin{chapteropening}
\textbf{By the end of this chapter:} \emph{I can say condolence, a tribute, a death anniversary, a cemetery and to offer condolences.}

Complete the last lesson of chapter 392: I can say condolence, a tribute, a death anniversary, a cemetery and to offer condolences.
\end{chapteropening}

\section[\texorpdfstring{\ml{അനുശോചനം} (anuśōcanaṁ)}{അനുശോചനം (anuśōcanaṁ)}]{\ml{അനുശോചനം} (anuśōcanaṁ) — condolence}

\label{lesson:ML-C392-anusocanam}



\begin{quote}
Before the new one: say the Malayalam for to cremate, then the Malayalam for to bury.
\end{quote}

\subsection*{You'll want to know: \ml{അനുശോചനം}}
\textbf{\ml{അനുശോചനം}} — \emph{anuśōcanaṁ} — \textquotedblleft{}condolence\textquotedblright{}.

\begin{quote}
\ml{മാനേജർ} \ml{കുടുംബത്തിന്} \ml{അനുശോചനം} \ml{അറിയിച്ചു}.
\end{quote}

\begin{quote}
\emph{mānējar kuṭumbattinŭ anuśōcanaṁ aṟiyiccu.}
\end{quote}

\begin{quote}
The manager conveyed his condolences to the family.
\end{quote}

Condolences are conveyed: \ml{അനുശോചനം} \ml{അറിയിക്കുക}.

\begin{cousinweb}
\textbf{\ml{അനുശോചനം}} (\emph{anuśōcanaṁ}) is borrowed from Sanskrit \emph{anuśocana}, \textquotedblleft{}sorrow, repentance\textquotedblright{} (Monier-Williams), from \emph{anu-śuc}, \textquotedblleft{}to grieve over\textquotedblright{}.
\end{cousinweb}

\begin{grammarlens}[title={What you've built}]
One more word for telling what happened.
\end{grammarlens}

\subsection*{Your turn}
\begin{itemize}
\raggedright
  \item \emph{Say it:} \emph{anuśōcanaṁ}
  \item \emph{Say it:} \emph{anuśōcanaṁ}, once more
  \item \emph{Say it:} \emph{mānējar kuṭumbattinŭ anuśōcanaṁ aṟiyiccu.} — The manager conveyed his condolences to the family.
  \item \emph{Recall:} say the Malayalam for to cremate, then the Malayalam for to bury, then say \emph{anuśōcanaṁ} again
\end{itemize}

\begin{tcolorbox}[breakable,colback=teal!4,colframe=teal!35!black,title={Before you move on}]
What does \ml{അനുശോചനം} mean? (\textquotedblleft{}Condolence\textquotedblright{}.) Say it once more.
\end{tcolorbox}
\section[\texorpdfstring{\ml{ആദരാഞ്ജലി} (ādarāñjali)}{ആദരാഞ്ജലി (ādarāñjali)}]{\ml{ആദരാഞ്ജലി} (ādarāñjali) — a tribute; homage}

\label{lesson:ML-C392-adaranjali}



\begin{quote}
Before the new one: say the Malayalam for to bury, then the Malayalam for condolence.
\end{quote}

\subsection*{You'll want to know: \ml{ആദരാഞ്ജലി}}
\textbf{\ml{ആദരാഞ്ജലി}} — \emph{ādarāñjali} — \textquotedblleft{}a tribute; homage\textquotedblright{}.

\begin{quote}
\ml{നാട്ടുകാർ} \ml{കവിക്ക്} \ml{ആദരാഞ്ജലി} \ml{അർപ്പിച്ചു}.
\end{quote}

\begin{quote}
\emph{nāṭṭukār kavikkŭ ādarāñjali arppiccu.}
\end{quote}

\begin{quote}
The local people paid tribute to the poet.
\end{quote}

It is offered: \ml{ആദരാഞ്ജലി} \ml{അർപ്പിക്കുക}.

\begin{cousinweb}
\textbf{\ml{ആദരാഞ്ജലി}} (\emph{ādarāñjali}) is built from two Sanskrit words, \emph{ādara}, \textquotedblleft{}respect\textquotedblright{} (CDIAL 1142), and \emph{añjali}, \textquotedblleft{}the hollowed hands placed together\textquotedblright{} (CDIAL 171): hands joined in respect.
\end{cousinweb}

\begin{grammarlens}[title={What you've built}]
One more word for telling what happened.
\end{grammarlens}

\subsection*{Your turn}
\begin{itemize}
\raggedright
  \item \emph{Say it:} \emph{ādarāñjali}
  \item \emph{Say it:} \emph{ādarāñjali}, once more
  \item \emph{Say it:} \emph{nāṭṭukār kavikkŭ ādarāñjali arppiccu.} — The local people paid tribute to the poet.
  \item \emph{Recall:} say the Malayalam for to bury, then the Malayalam for condolence, then say \emph{ādarāñjali} again
\end{itemize}

\begin{tcolorbox}[breakable,colback=teal!4,colframe=teal!35!black,title={Before you move on}]
What does \ml{ആദരാഞ്ജലി} mean? (\textquotedblleft{}A tribute; homage\textquotedblright{}.) Say it once more.
\end{tcolorbox}
\section[\texorpdfstring{\ml{ചരമവാർഷികം} (caramavārṣikaṁ)}{ചരമവാർഷികം (caramavārṣikaṁ)}]{\ml{ചരമവാർഷികം} (caramavārṣikaṁ) — a death anniversary}

\label{lesson:ML-C392-caramavarsikam}



\begin{quote}
Before the new one: say the Malayalam for condolence, then the Malayalam for a tribute.
\end{quote}

\subsection*{You'll want to know: \ml{ചരമവാർഷികം}}
\textbf{\ml{ചരമവാർഷികം}} — \emph{caramavārṣikaṁ} — \textquotedblleft{}a death anniversary\textquotedblright{}.

\begin{quote}
\ml{ഇന്ന്} \ml{മുത്തച്ഛന്റെ} \ml{ചരമവാർഷികം} \ml{ആണ്}.
\end{quote}

\begin{quote}
\emph{innŭ muttacchanṟe caramavārṣikaṁ āṇŭ.}
\end{quote}

\begin{quote}
Today is grandfather's death anniversary.
\end{quote}

\ml{ചരമം}, a death, + \ml{വാർഷികം}, an anniversary.

\begin{cousinweb}
\textbf{\ml{ചരമവാർഷികം}} (\emph{caramavārṣikaṁ}) is built from \ml{ചരമം} (\emph{caramaṁ}), \textquotedblleft{}a death\textquotedblright{}, from Sanskrit \emph{carama}, \textquotedblleft{}last\textquotedblright{} (Monier-Williams), and \ml{വാർഷികം} (\emph{vārṣikaṁ}), \textquotedblleft{}an anniversary\textquotedblright{}, from \emph{vārṣika}, \textquotedblleft{}lasting a year, annual\textquotedblright{} (CDIAL 11571).
\end{cousinweb}

\begin{grammarlens}[title={What you've built}]
One more word for telling what happened.
\end{grammarlens}

\subsection*{Your turn}
\begin{itemize}
\raggedright
  \item \emph{Say it:} \emph{caramavārṣikaṁ}
  \item \emph{Say it:} \emph{caramavārṣikaṁ}, once more
  \item \emph{Say it:} \emph{innŭ muttacchanṟe caramavārṣikaṁ āṇŭ.} — Today is grandfather's death anniversary.
  \item \emph{Recall:} say the Malayalam for condolence, then the Malayalam for a tribute, then say \emph{caramavārṣikaṁ} again
\end{itemize}

\begin{tcolorbox}[breakable,colback=teal!4,colframe=teal!35!black,title={Before you move on}]
What does \ml{ചരമവാർഷികം} mean? (\textquotedblleft{}A death anniversary\textquotedblright{}.) Say it once more.
\end{tcolorbox}
\section[\texorpdfstring{\ml{സെമിത്തേരി} (semittēri)}{സെമിത്തേരി (semittēri)}]{\ml{സെമിത്തേരി} (semittēri) — a cemetery}

\label{lesson:ML-C392-semitteri}



\begin{quote}
Before the new one: say the Malayalam for a tribute, then the Malayalam for a death anniversary.
\end{quote}

\subsection*{You'll want to know: \ml{സെമിത്തേരി}}
\textbf{\ml{സെമിത്തേരി}} — \emph{semittēri} — \textquotedblleft{}a cemetery\textquotedblright{}.

\begin{quote}
\ml{പള്ളിയുടെ} \ml{പിന്നിൽ} \ml{ഒരു} \ml{പഴയ} \ml{സെമിത്തേരി} \ml{ഉണ്ട്}.
\end{quote}

\begin{quote}
\emph{paḷḷiyuṭe pinnil oru paḻaya semittēri uṇṭŭ.}
\end{quote}

\begin{quote}
There is an old cemetery behind the church.
\end{quote}

The burial ground of a church.

\begin{cousinweb}
\textbf{\ml{സെമിത്തേരി}} (\emph{semittēri}) is usually traced to Portuguese \emph{cemitério}, \textquotedblleft{}a cemetery\textquotedblright{}: a church word that came with the Portuguese to the coast.
\end{cousinweb}

\begin{grammarlens}[title={What you've built}]
One more word for telling what happened.
\end{grammarlens}

\subsection*{Your turn}
\begin{itemize}
\raggedright
  \item \emph{Say it:} \emph{semittēri}
  \item \emph{Say it:} \emph{semittēri}, once more
  \item \emph{Say it:} \emph{paḷḷiyuṭe pinnil oru paḻaya semittēri uṇṭŭ.} — There is an old cemetery behind the church.
  \item \emph{Recall:} say the Malayalam for a tribute, then the Malayalam for a death anniversary, then say \emph{semittēri} again
\end{itemize}

\begin{tcolorbox}[breakable,colback=teal!4,colframe=teal!35!black,title={Before you move on}]
What does \ml{സെമിത്തേരി} mean? (\textquotedblleft{}A cemetery\textquotedblright{}.) Say it once more.
\end{tcolorbox}
\section[\texorpdfstring{\ml{അനുശോചിക്കുക} (anuśōcikkuka)}{അനുശോചിക്കുക (anuśōcikkuka)}]{\ml{അനുശോചിക്കുക} (anuśōcikkuka) — to offer condolences}

\label{lesson:ML-C392-anusocikkuka}



\begin{quote}
Before the new one: say the Malayalam for a death anniversary, then the Malayalam for a cemetery.
\end{quote}

\subsection*{You'll want to know: \ml{അനുശോചിക്കുക}}
\textbf{\ml{അനുശോചിക്കുക}} — \emph{anuśōcikkuka} — \textquotedblleft{}to offer condolences\textquotedblright{}.

\begin{quote}
\ml{അയൽക്കാർ} \ml{ആ} \ml{മരണത്തിൽ} \ml{അനുശോചിച്ചു}.
\end{quote}

\begin{quote}
\emph{ayalkkār ā maraṇattil anuśōciccu.}
\end{quote}

\begin{quote}
The neighbours offered condolences on that death.
\end{quote}

Past: \ml{അനുശോചിച്ചു}, \emph{anuśōciccu}.

The death takes -\ml{ഇൽ}: \ml{മരണത്തിൽ} \ml{അനുശോചിച്ചു}.

\begin{cousinweb}
\textbf{\ml{അനുശോചിക്കുക}} (\emph{anuśōcikkuka}) is built on \ml{അനുശോചനം} (\emph{anuśōcanaṁ}), from Sanskrit \emph{anuśocana}, \textquotedblleft{}sorrow\textquotedblright{} (Monier-Williams), with the Malayalam verb ending -\ml{ിക്കുക}.
\end{cousinweb}

\begin{grammarlens}[title={What you've built}]
One more word for telling what happened.
\end{grammarlens}

\subsection*{Your turn}
\begin{itemize}
\raggedright
  \item \emph{Say it:} \emph{anuśōcikkuka}
  \item \emph{Say it:} \emph{anuśōcikkuka}, once more
  \item \emph{Say it:} \emph{ayalkkār ā maraṇattil anuśōciccu.} — The neighbours offered condolences on that death.
  \item \emph{Recall:} say the Malayalam for a death anniversary, then the Malayalam for a cemetery, then say \emph{anuśōcikkuka} again
\end{itemize}

\begin{tcolorbox}[breakable,colback=teal!4,colframe=teal!35!black,title={Before you move on}]
What does \ml{അനുശോചിക്കുക} mean? (\textquotedblleft{}To offer condolences\textquotedblright{}.) Say it once more.
\end{tcolorbox}
